\chapter{Federated Glucose Forecasting: Benchmark and Live Deployment}
\label{ch:glucose}
% Target: ~6,000 words. From Paper 2.

\section{Introduction}
\outline{Almost all federated glucose work is simulated. Methods are rarely compared on one task.}

\section{Data}
\outline{Four T1D cohorts as federated clients. Three unseen cohorts (ReplaceBG, BrisT1D, Flair) as out-of-distribution tests. Ethics and data access per cohort.}

\section{Methods}
\outline{One encoder-decoder transformer. By-patient splits, five seeds. Strategies compared: local training, FedAvg, FedProx, MLDG, APFL, Ditto, plus a pooled reference.}

\section{The live federation}
\outline{Four UK universities (UCL, Manchester, Newcastle, Oxford), each holding only its own data. What the tooling is and why it is light enough for clinics.}

\section{Results}
\outline{Federation cost almost nothing: within 0.09 mg/dL RMSE of pooled training, far below CGM sensor error. Single-site training is a gamble. One local model matched federation out of distribution, another collapsed, and in-domain accuracy gave no warning. A new site beat training from scratch after finetuning on 9 to 17 patients. Differences between FL strategies fell within seed noise. The dominant effect is federation itself.}

\section{Discussion and limitations}
\outline{From the paper. Metric choice, seed count, cohort coverage, security of the channel.}
